\title{Controlling Text-to-Image Diffusion by Orthogonal Finetuning}
The field of text-to-image synthesis has advanced rapidly~\cite{saharia2022photorealistic,ramesh2022hierarchical,rombach2022high,gu2022vector,nichol2022glide}, fueled by significant innovations in diffusion-based generative models~\cite{ho2020denoising,song2021score,song2021denoising,dhariwal2021diffusion} and advancements in vision-language joint representations~\cite{su2020vlbert,lu2019vilbert,radford2021learning,wang2021simvlm,li2022blip,li2023blip,alayrac2022flamingo,schuhmann2022laion}. This process can also be interpreted as modifying the standard coordinate basis for all neurons within a layer. To empirically verify our reasoning, we conduct a controlled experiment modeled after the ControlNet~\cite{zhang2023adding} setup; the results, displayed in Figure~\ref{fig:ortho}, show that our typical OFT yields consistent and precise control at the end of training (epoch 20). This metric operates by extracting the control signal from the output image and measuring its agreement with the initial input control signal. Reviewing Figure~\ref{fig:dreambooth_final}, it is evident that OFT and COFT both offer robust semantic retention of the subject, whereas LoRA frequently struggles to maintain the subject’s identity (\emph{e.g.}, the bowl instance, where LoRA entirely loses subject identity). To ensure an equitable evaluation, each sample is produced using the same finetuned model across all approaches, and we refrain from optimizing the text prompt individually for each output. The rationale behind employing orthogonal transformations for neuron finetuning is partly inspired by principles from the 2D Fourier transform, a technique that decomposes images into their magnitude and phase spectra. Broadly, LoRA's maximum attainable performance remains modest, even after rigorous hyperparameter tuning. \textbf{Controllable generation}. For subject-driven generation, we adopt the methodology described in DreamBooth~\cite{ruiz2023dreambooth}, and for controllable generation tasks, we largely follow protocols established by ControlNet~\cite{zhang2023adding} and T2I-Adapter~\cite{mou2023t2i}. Every method undergoes finetuning on Stable Diffusion~(SD) v1.5 for 20 epochs per dataset. For both DreamBooth and LoRA, we adhere closely to the procedures and optimal hyperparameter choices recommended in their respective original publications. \subsection{Constrained Orthogonal Finetuning}

The expressivity of the standard OFT method can be further curtailed by restricting the learned transformation to remain in close proximity to the initial pretrained model. Adapting large pretrained models to downstream applications by finetuning has become standard practice~\cite{devlin2018bert,brown2020language,he2022masked}. DreamBooth~\cite{ruiz2023dreambooth} and LoRA~\cite{hulora2022} are employed as baseline methods. Additionally, OFT demonstrates higher sample efficiency, achieving satisfactory results with substantially fewer images and fewer training epochs. Besides diffusion techniques, text-to-image synthesis has been explored via generative adversarial networks~\cite{reed2016generative,zhang2017stackgan,xu2018attngan,li2019controllable} and autoregressive frameworks~\cite{ramesh2021zero,ding2021cogview,wu2022nuwa,yuscaling}. We also present the results of a human evaluation in Appendix~\ref{app:human}, providing further support for the efficacy of OFT. \subsection{General Framework}

In standard finetuning, gradient descent is performed using a modest learning rate to adjust either the whole model or select layers. In the C2I task, OFT and COFT convincingly synthesize semantically coherent images from rough Canny edge inputs, a scenario where both T2I-Adapter and LoRA are markedly less effective. To illustrate, reconstructing an image using its original phase spectrum and a randomly generated magnitude spectrum largely retains the semantics of the source image. LoRA introduces a low-rank adaptation matrix to ensure that the parameters inherited from pretraining are not substantially altered. However, these straightforward strategies do not guarantee that the model's original generative strengths are maintained. Put another way, classic finetuning operates under an implicit objective of minimizing $\thickmuskip=2mu \medmuskip=2mu \| \bm{M} - \bm{M}^0\|$, with $\bm{M}$ representing the updated (finetuned) weights and $\bm{M}^0$ the original, pretrained weights. Image-to-image translation is essentially a problem of controllable content generation, and traditional approaches predominantly employ conditional generative adversarial networks~\cite{isola2017image,zhu2017toward,wang2018high,park2019semantic,choi2018stargan} for this purpose. To make this more concrete, consider a weight matrix of size $\thickmuskip=2mu \medmuskip=2mu 128\times 128$: for $\thickmuskip=2mu \medmuskip=2mu r' = 8$, LoRA results in $2{,}048$ parameters; for the same $r = 8$, OFT contains $960$ trainable parameters (here, block sharing is not used). \textbf{Quantitative comparison}. For instance, LoRA requires 20 epochs to match the accuracy OFT attains by epoch 8. Initial experimental results regarding convolution layer finetuning with OFT can be found in Appendix~\ref{app:convolution}. \textbf{Balancing flexibility and regularity in finetuning}. The block-diagonal configuration of $\bm{R}$ enables compelling interpretations specific to convolutional layers—something LoRA lacks. \section{Intriguing Insights and Discussions}

\textbf{OFT’s versatility across architectures}. This approach works by predicting the control signal present in the synthesized image and comparing it to the initial control signal. For a general orthogonal matrix of dimensions $\thickmuskip=2mu \medmuskip=2mu d\times d$, there are $\thickmuskip=2mu \medmuskip=2mu d(d-1)/2$ parameters to learn, corresponding to a computational complexity of $\mathcal{O}(d^2)$. Furthermore, the literature lacks rigorous principles for developing robust finetuning techniques and for determining crucial hyperparameters, such as training duration. Additional experimental findings are detailed in Appendix~\ref{app:settings},\ref{app:coft_oft},\ref{app:more_qualitative},\ref{app:failure}. An effective approach to finetuning should satisfy two primary criteria: (1) \emph{training efficiency}, characterized by minimizing the number of parameters adjusted and reducing training time, and (2) \emph{generalizability preservation}, ensuring the model retains high-fidelity and diversity in its outputs. Both GLIDE~\cite{nichol2022glide} and Imagen~\cite{saharia2022photorealistic} operate their diffusion processes in pixel space; the former employs a joint text encoder-diffusion prior trained on paired data, whereas the latter leverages a fixed, pretrained text encoder. To make the finetuning process more restrictive, LoRA enforces an explicit low-rank structure on the updates: $\thickmuskip=2mu \medmuskip=2mu \text{rank}(\bm{M} - \bm{M}^0)=r'$, where $r'$ is a user-defined small integer. This characteristic allows both OFT and COFT to be finetuned with a generously high iteration count without the need for fine-tuning. This extension, re-scaled OFT, further broadens the adaptability of OFT with only a negligible increase in parameter count. This observation points to the pivotal role of adjusting neuron orientations for semantic manipulation in generated images, aligning with objectives in both subject-driven and controllable generation. \item We deploy OFT across two application domains: subject-driven generation and controllable generation. According to the results in Table~\ref{table:control_gen}, both OFT and COFT deliver notably superior and more precise control compared to alternative techniques. Furthermore, the performance of OFT remains stable regardless of the number of iterations, and thus, OFT retains effectiveness even with many iterations—a property not observed for DreamBooth or LoRA. To assess the convergence rate of ControlNet, T2I-Adapter, LoRA, and COFT on the L2F task, we carry out both quantitative and qualitative analyses. \textbf{Qualitative comparison}. This is a typographical error; the original implementation (used for Appendix~\ref{app:rescaled}) is correct.} where $\thickmuskip=2mu \medmuskip=2mu \bm{D}=\text{diag}(s_1,\cdots,s_n)\in\mathbb{R}^{n\times n}$ is a diagonal matrix whose trainable diagonal entries $s_1,\dots,s_n$ are restricted to be positive. While this intuition sometimes aligns with practice, we contend that relying solely on the Euclidean norm fails to reflect the full scope of semantic retention. Similarly,~\cite{nitzan2022mystyle} learns a tailored generative face prior from an individual's image set. For an $r$-block diagonal orthogonal matrix, the parameter count becomes $\thickmuskip=2mu \medmuskip=2mu d(d/r-1)/2$, resulting in a complexity of $\thickmuskip=2mu \medmuskip=2mu \mathcal{O}(d^2/r)$. \end{itemize}

At their core, both tasks address the challenge of how to adapt text-to-image diffusion models via finetuning, while maintaining the generative capabilities achieved during pretraining. The smooth and reliable convergence behavior characteristic of OFT not only simplifies practical deployment but also makes it easier to tune the method’s hyperparameters effectively. \textbf{Connection to LoRA}. SphereNet~\cite{liu2017deep} demonstrates that normalizing all neural activations onto a unit hypersphere during training retains model capacity and even enhances generalization, suggesting that the orientation of neural representations encapsulates crucial data information. This low learning rate inherently constrains the adjustments, encouraging only minor alterations from the pretrained parameter values. To ensure a direct comparison with LoRA, our OFT finetuning experiments are likewise limited to the attention weights. Each method is evaluated under its optimal hyperparameter configuration. Notably, adapter-based strategies (\emph{e.g.}, T2I-Adapter and ControlNet) demonstrate slower convergence and poorer ultimate performance. In addition, OFT readily extends to convolutional architectures. The optimal value for Eq.~\eqref{eq:oft_principle} is trivially zero and is reached whenever $\bm{W}$ and $\bm{W}^0$ are related through an orthogonal transformation, i.e., $\thickmuskip=2mu \medmuskip=2mu \bm{W}=\bm{R}\bm{W}^0$ where $\thickmuskip=2mu \medmuskip=2mu \bm{R}\in\mathbb{R}^{d\times d}$ is an orthogonal matrix (signifying a rotation if $\det \bm{R}=1$, or a reflection if $\det \bm{R}=-1$). We observe that maintaining this property is essential for retaining the semantic fidelity in the generative process of these models. To ensure impartiality, generated outputs are randomly selected for all compared methods. \section{Orthogonal Finetuning}

\subsection{Why Does Orthogonal Transformation Make Sense?}

We begin by examining the rationale for employing orthogonal transformations in the finetuning of text-to-image diffusion models. \subsection{Re-scaled Orthogonal Finetuning}

We introduce a straightforward modification to standard OFT by learning an additional scaling factor per neuron. Despite their success, these models can struggle with ambiguous text instructions and often lack the capacity for precise, detailed manipulation of image attributes. An outstanding question, however, remains: how can these advanced models be steered or adapted for a variety of downstream applications in a controlled manner? Specifically, for the C2I task, we calculate the IoU and F1 score, while for S2I, we determine both the mean IoU and mean/overall accuracy. We further perform a qualitative evaluation of OFT and COFT alongside established state-of-the-art approaches such as ControlNet, T2I-Adapter, and LoRA. This approach mathematically ensures the preservation of pairwise angular relationships among neurons in a layer, thereby promoting greater stability during finetuning. This methodology leverages the fact that randomly initialized networks have, in expectation, low hyperspherical energy, and the objective in~\cite{liu2021orthogonal} is to sustain low hyperspherical energy throughout training, which correlates with improved generalization performance~\cite{liu2018learning,lin2020regularizing}. \textbf{Why not minimize hyperspherical energy}. This is accomplished by initializing the orthogonal matrix $\bm{R}$ to the identity, which guarantees the pretrained weights as the starting point for finetuning. For our qualitative comparison, we restrict the discussion to OFT, COFT, and ControlNet, as ControlNet yields the closest landmark error to ours. The results in~\cite{liu2021orthogonal} suggest orthogonal transformations provide ample flexibility to develop neural networks capable of strong classification generalization. Adopting the evaluation protocol in \cite{ruiz2023dreambooth}, we quantitatively measure subject fidelity (DINO~\cite{caron2021emerging}, CLIP-I~\cite{radford2021learning}), text prompt fidelity (CLIP-T~\cite{radford2021learning}), and sample diversity (LPIPS~\cite{zhang2018unreasonable}). It is possible to further decrease the number of trainable parameters by sharing the blocks, i.e., letting $\thickmuskip=2mu \medmuskip=2mu \bm{R}_i = \bm{R}_j$ for all $i\neq j$, which leads to a parameter complexity of $\mathcal{O}(d^2/r^2)$. With this perspective, a natural approach is to employ a shared orthogonal transformation at each network layer, ensuring the hyperspherical energy remains invariant. When juxtaposed with ControlNet, LoRA achieves higher scores for the S2I task but underperforms in C2I and L2F. Interestingly, OFT, despite its straightforward formulation, effectively strikes a balance by maintaining the angular relationships among neuron pairs. Existing solutions aiming to avoid unwanted modifications to the main subject, such as~\cite{avrahami2022blended,nichol2022glide}, rely on user-provided masks as conditional inputs. While greater adaptability could theoretically be achieved through more general angular modifications—such as rotating individual neurons independently—we observe that orthogonal transformations strike an optimal balance between adaptability and structural discipline. Further,~\cite{hertz2022prompt} is capable of performing targeted local and holistic edits without the need for external masks. To enhance stability during finetuning, we additionally introduce Constrained Orthogonal Finetuning (COFT), which incorporates a supplementary radius restriction on the hypersphere. Stable Diffusion integrates VQ-GAN~\cite{esser2021taming} to construct a latent visual vocabulary, while DALL-E2 utilizes CLIP~\cite{radford2021learning} for a shared embedding space between text and images. While this implicit regularization is conceptually reasonable, it may still allow too much flexibility, especially when adapting large-scale models. For all methods under comparison, the model exhibiting the highest validation CLIP metrics is chosen. Key distinctions between our Orthogonal Finetuning (OFT) approach and the method from~\cite{liu2021orthogonal} are twofold. Our working assumption is that an optimally finetuned model should exhibit a hyperspherical energy very close to that of its pretrained counterpart. The introduction of ControlNet~\cite{zhang2023adding} enables explicit control over pretrained diffusion models by leveraging finetuning strategies and the integration of auxiliary control signals, thereby enhancing controllability. To better understand the properties of the Cayley transform, one may approximate $\thickmuskip=2mu \medmuskip=2mu \bm{R}=(\bm{I}+\bm{Q})(I-\bm{Q})^{-1}$ by

series converges under the operator norm). Regarding the first question, we build on the empirical findings from \cite{liu2018decoupled,chen2020angular}, which indicate that the angular difference between features effectively reflects semantic disparities. Relative to ControlNet, OFT and COFT produce images with enhanced detail, more natural geometric composition, and accomplish this with a significantly reduced parameter count. Importantly, irrespective of these parameter reduction techniques, the overall matrix $\bm{R}$ retains its orthogonality, thus fully maintaining hyperspherical energy preservation. In comparison, OFT employs an orthogonal (and thus full-rank) transformation applied to the pretrained weight matrix, thereby preserving the critical pretrained knowledge by constraining the transformation. While the orthogonality property is easily enforced using Cayley parameterization as detailed in Section~\ref{sect:eff_ortho}, integrating the $\epsilon$-closeness constraint into a Cayley-parametrized orthogonal matrix is more challenging. \textbf{Qualitative comparison}. Importantly, our methods exhibit limited sensitivity to the number of finetuning iterations, greatly reducing the challenge of determining the ideal iteration count for various subjects. Empirically, this constraint has negligible impact on practical performance. Additionally, OFT and COFT demonstrate significantly better adherence to text prompts, in contrast to DreamBooth, which, although successful at preserving subject identity, often fails to generate outputs consistent with the provided textual descriptions (\emph{e.g.}, the first row for the bowl example). \end{itemize}

\section{Related Work}

\textbf{Text-to-image diffusion models}. For baselines, we consider ControlNet~\cite{zhang2023adding}, T2I-Adapter~\cite{mou2023t2i}, and LoRA~\cite{hulora2022}. \subsection{Efficient Orthogonal Parameterization}\label{sect:eff_ortho}

Traditional procedures for orthogonalization, such as the Gram-Schmidt process, while differentiable, tend to be computationally prohibitive in large-scale scenarios~\cite{liu2021orthogonal}. Second, OFT's structure allows for promising exploration in compositionality: the orthogonal matrices generated by distinct OFT finetuning tasks can be composed through matrix multiplication and still retain orthogonality. Additionally, considering that minimizing the Euclidean distance between the pretrained and finetuned models favors retention of the pretrained model’s capabilities, we further develop a variant termed Constrained Orthogonal Finetuning~(COFT), which enforces the finetuned model to reside within a fixed-radius hypersphere centered at the positions of the pretrained neurons. \section{Experiments and Results}

\textbf{Experimental setup}. Supporting this, \cite{liu2021orthogonal} establishes that orthogonal transformations possess sufficient representational strength for neural network learning. Our investigation reveals a fundamental trade-off between flexibility and regularity during model finetuning. To address this, we introduce a constraint that maintains the angle between every neuron pair, motivated by the premise that these inter-neuron angles are fundamental to the representation of knowledge in neural networks. Consistent with the frameworks outlined in \cite{zhang2023adding,mou2023t2i}, we utilize OFT to finetune pretrained diffusion models, achieving improved controllability while requiring fewer data samples and a smaller number of trainable parameters. Consequently, the constraint $\thickmuskip=2mu \medmuskip=2mu\|\bm{R}-\bm{I}\|\leq\epsilon$ may be incorporated within the Cayley transform framework, translating into the condition $\thickmuskip=2mu \medmuskip=2mu \|\bm{Q}-\bm{0}\|\leq\epsilon'$ where $\bm{0}$ is the zero matrix and $\epsilon'$ serves as a distinct error parameter from $\epsilon$. However, this approach does not necessarily guarantee minimal deviation in hyperspherical energy. Notably, the cosine-based (c) approach is not employed in the training phase—training exclusively uses the inner product. Supplementary results and methodological specifics can be found in Appendix~\ref{app:settings}. Just as LoRA modulates trainable parameters through its rank parameter $r'$, OFT provides a diagonal block size $r$ to control parameter efficiency. OFT adopts the overall methodology of DreamBooth but substitutes standard finetuning with a parameter update rule based on orthogonal transformations instead of employing a mere low learning rate. Expanded results are provided in Appendix~\ref{app:more_qualitative}, \ref{app:more_control_task}, and \ref{app:failure}. The outcomes are depicted in Figure~\ref{fig:teaser} and Figure~\ref{fig:db_convergence}. Notably, the hyperparameter $\epsilon$ modulates the range of adaptation during finetuning: increasing $\epsilon$ allows the COFT model’s behavior to approach that of OFT, while reducing $\epsilon$ causes COFT's output to more closely align with the unadapted pretrained diffusion model. OFT is particularly related to LoRA~\cite{hulora2022}, which assumes the finetuning updates are of low-rank form and are additive in nature. CLIP-I quantifies the mean pairwise cosine similarity between CLIP embeddings of the generated and actual images, while DINO follows a similar procedure using ViT S/16 DINO embeddings. More recent advances have incorporated diffusion models into image-to-image translation pipelines~\cite{saharia2022palette,voynov2022sketch,wang2022pretraining}. Hyperspherical energy is quantified by summing the hyperspherical similarities (\emph{e.g.}, cosine similarities) between all neuron pairs within a layer, thus reflecting the degree of neuron dispersion on the unit hypersphere~\cite{liu2021learning}. A straightforward strategy might involve incorporating a regularization term to maintain the hyperspherical energy constant throughout finetuning. OFT's forward computation can be expressed as $\thickmuskip=2mu \medmuskip=2mu \bm{z}=(\bm{R}\bm{W}^0)^\top\bm{x}=(\bm{W}^0+(\bm{R}-\bm{I})\bm{W}^0)^\top\bm{x}$, where $\thickmuskip=2mu \medmuskip=2mu (\bm{R}-\bm{I})\bm{W}^0$ plays a role comparable to LoRA's modification via a low-rank update. Crucially, OFT maintains computational efficiency by not adding any additional inference burden, thus making it highly suitable for practical deployment. In the S2I task, LoRA is unable to adhere to the provided segmentation maps, whereas ControlNet, OFT, and COFT reliably guide the generative process to respect input conditions. \textbf{Finetuning stability and convergence}. The superior convergence speed of OFT is further substantiated by Figure~\ref{fig:teaser}, where the example on the right illustrates that OFT achieves effective convergence using just 5\% of the ADE20K training set, outperforming both ControlNet and LoRA in terms of data efficiency. Distinct from current approaches, OFT mathematically guarantees the preservation of hyperspherical energy, a metric encapsulating the inter-neuron pairwise relationships on the unit hypersphere. Unlike the approach in \cite{liu2021orthogonal}, we do not seek to optimize for minimum hyperspherical energy. Further information about these consistency evaluation protocols can be found in Appendix~\ref{app:settings}. For L2F, the evaluation is based on the average $\ell_2$ distance between the input and predicted landmarks. This is the

\begin{equation}\label{eq:oft_general}
    \footnotesize
    \bm{z}=\bm{W}^\top\bm{x}=(\bm{R}\cdot\bm{W}^0)^\top\bm{x},~~~\text{s.t.}~\bm{R}^\top\bm{R}=\bm{R}\bm{R}^\top=\bm{I}
\end{equation}

In this expression, $\bm{W}$ is the weight matrix updated via OFT-finetuning and $\bm{I}$ refers to the identity matrix. We begin by assessing both the stability and speed of convergence during finetuning for DreamBooth, LoRA, OFT, and COFT. Present methods generally assume that maintaining a small Euclidean distance between the parameters of the finetuned and pretrained models suffices for preserving these abilities, leading to the prevalent use of small learning rates during finetuning. If both $r$ and $r'$ are set proportionally to the neuron dimension $d$ (for example, $\thickmuskip=2mu \medmuskip=2mu r = r' = \alpha d$, with $\thickmuskip=2mu \medmuskip=2mu 0<\alpha\leq 1$), LoRA exhibits a parameter complexity of $\mathcal{O}(d^2+dn)$, whereas OFT achieves $\mathcal{O}(d)$. The resulting output $\thickmuskip=2mu \medmuskip=2mu \bm{z}\in\mathbb{R}^n$ is determined via $\thickmuskip=2mu \medmuskip=2mu \bm{z}=\bm{W}^\top\bm{x}$, with $\thickmuskip=2mu \medmuskip=2mu \bm{x}\in\mathbb{R}^d$ serving as the input. Concretely, our modified forward computation is given by: $\thickmuskip=2mu \medmuskip=2mu\bm{z}=(\bm{R}\bm{W}^0\bm{D})^\top\bm{x}$\footnote{\textbf{Errata}: In the NeurIPS camera-ready version, the formula for the re-scaled OFT forward pass was erroneously given as $\thickmuskip=2mu \medmuskip=2mu\bm{z}=(\bm{D}\bm{R}\bm{W}^0)^\top\bm{x}$. To bolster training stability, we also develop a constrained version that restricts the angular displacement from the original pretrained parameters. Whether this composition preserves the knowledge acquired for each downstream task is an intriguing topic for further research. The data demonstrate that both COFT and OFT reach lower error rates notably faster than competing approaches. To offer a fuller picture of these qualitative improvements, we include additional visual examples for each control task in Appendix~\ref{app:control_exp}, and further show that OFT generalizes well to a broader spectrum of control applications—such as dense pose-to-body, sketch-to-image, and depth-to-image translation—in Appendix~\ref{app:more_control_task}. This efficient technique introduces a minor constraint: the Cayley parameterization only generates orthogonal matrices with determinant $1$, limiting the output to the special orthogonal group. The CLIP-T metric corresponds to the average cosine similarity between the text prompt’s CLIP embedding and that of the generated images. To clarify this rationale, we decompose the overarching question into two components: (1) what motivates finetuning of the angular (directional) properties of neuron weights, and (2) what makes orthogonal transformation a suitable method for adjusting these angles. Conventionally, finetuning is performed either by applying small learning rates to modify existing neuron weights (\emph{e.g.}, \cite{ruiz2023dreambooth}) or by augmenting the network with compact modules containing re-parameterized weights (\emph{e.g.}, \cite{hulora2022,zhang2023adding}). Figure~\ref{fig:face_convergence} presents the evolution of face landmark error as a function of finetuning epochs for each method. In our main experiments, we use the original OFT formulation in order to highlight the orthogonal transformation itself, though we generally observe that re-scaled OFT achieves superior performance (refer to Appendix~\ref{app:rescaled} for further analysis). For the L2F task, our approach consistently achieves precise pose alignment, even in cases featuring complex facial orientations. Through extensive empirical evaluations, we reveal notable gains over state-of-the-art methods in terms of sample quality, training convergence speed, and stability during finetuning. All models are optimized using the same loss formulation as specified in DreamBooth. These qualitative findings underscore that the OFT framework achieves superior controllability alongside reliable subject retention. DreamBooth~\cite{ruiz2023dreambooth}, using a custom token and a handful of subject images, adapts a text-to-image diffusion model through a combination of reconstruction and class-preserving losses. Consequently, the inference throughput matches that of the base pretrained network. \end{document} To further illustrate the significance of neuron angles, we carry out a simple experiment depicted in Figure~\ref{fig:toy_ae}: we train a conventional convolutional autoencoder with flower image inputs. Secondly, OFT introduces constraints to limit divergence from the pretrained network, leading to a constrained approach, whereas~\cite{liu2021orthogonal} operates without such restrictions. As illustrated by the randomly sampled images in Figure~\ref{fig:control_final}, both OFT and COFT demonstrate superior image realism and high fidelity, while simultaneously enabling precise control. \textbf{Subject-driven generation}. \section{Concluding Remarks and Open Problems}

Inspired by the critical role that angular relationships among neurons play in defining visual semantics, we introduce a straightforward yet powerful finetuning technique—orthogonal finetuning—for manipulating text-to-image diffusion models. When $\bm{R}-\bm{I}$ possesses low rank (which is often the case since $\bm{W}^0$ is typically full rank), OFT effectively enacts a low-rank perturbation to the weights. \item Unlike conventional finetuning techniques, OFT ensures that the hyperspherical energy of the model is mathematically maintained during training. Our metric confirms OFT’s superior controllability. Fundamentally, LoRA and OFT encapsulate two complementary principles for parameter-efficient learning: LoRA leverages low-rank adaptation for parameter reduction, while OFT utilizes a structured sparsity—namely, block-diagonal orthogonal transformations—to achieve similar ends. Conversely, OFT continues to improve as the parameter count increases, with no observable plateau in performance so far. Our experiments indicate that the OFT methodology leads to superior generative quality and faster convergence compared to leading existing approaches. Addressing the second question, the most straightforward approach to adjust neuron directions is to synchronously apply a rotation or reflection to all neurons within a layer, which inherently invokes an orthogonal transformation. Notably, the phase spectrum—representing the angular relationship between inputs and their corresponding bases—encapsulates much of the semantic content. By the 400th iteration, DreamBooth and both OFT variants demonstrate effective control, in contrast to LoRA, which does not maintain the subject’s identity. Furthermore, OFT may be interpreted as updating neurons while constraining them to move along a smooth hypersphere manifold, which is associated with a more favorable optimization geometry. Figure~\ref{fig:eps_coft} illustrates how varying $\epsilon$ influences COFT’s effectiveness. Conventional finetuning techniques offer the highest level of flexibility; however, they tend to compromise stability and are susceptible to model collapse. Results presented in Table~\ref{table:dreambooth_final} demonstrate that both COFT and OFT substantially outperform DreamBooth and LoRA on DINO and CLIP-I metrics, and either match or marginally surpass them regarding prompt fidelity and diversity. Relative to previous approaches, OFT offers enhanced control and more stable finetuning, all while utilizing a reduced set of finetuning parameters. Conversely, OFT and COFT consistently maintain robust and stable control over the outputs. Beyond 2000 iterations, DreamBooth suffers from mode collapse, and LoRA fails to render the yellow shirt, mistakenly generating yellow fur instead. The main paper investigates three demanding tasks in the context of controllable generation: Canny edge to image~(C2I) using the COCO dataset~\cite{lin2014microsoft}, segmentation map to image~(S2I) on the ADE20K dataset~\cite{zhou2017scene}, and landmark to face~(L2F) on the CelebA-HQ dataset~\cite{karras2018progressive,xia2021tedigan}. Figure~\ref{fig:block} offers a visual representation of OFT. To mitigate this, we express $\bm{R}$ in a block-diagonal structure consisting of $r$ individual blocks, resulting in the following formulation:

\begin{equation}
    \footnotesize
    \bm{R}=\text{diag}(\bm{R}_1,\bm{R}_2,\cdots,\bm{R}_r)=\begin{bmatrix}
    \bm{R}_1\in \text{O}(\frac{d}{r}) &  &\\
    &\ddots & \\
    & & \bm{R}_r\in \text{O}(\frac{d}{r})
    \end{bmatrix}\in \text{O}(d)
\end{equation}

where $\text{O}(d)$ represents the orthogonal group in $d$ dimensions, and $\thickmuskip=2mu \medmuskip=2mu \bm{R} \in \mathbb{R}^{d\times d}$ along with $\thickmuskip=2mu \medmuskip=2mu \bm{R}_i \in \mathbb{R}^{d/r\times d/r},~\forall i$ are both orthogonal matrices. As OFT does not rely on model-specific components, it is broadly compatible and can be integrated with any text-to-image diffusion architecture. However, affording extensive freedom in altering neuron directions risks compromising the generative capacity conferred during pretraining. Building upon this property, we introduce Orthogonal Finetuning~(OFT), a method designed to adapt large text-to-image diffusion models for downstream tasks while maintaining their hyperspherical energy. As such, employing more structurally aware measures to assess the divergence between the finetuned and pretrained models could substantially enhance both the consistency of finetuning and the preservation of original model performance. During training, the model produces its feature map using the conventional inner product, where $z$ represents the output of convolutional filter $\bm{w}$ applied to input $\bm{x}$ within a sliding window. In the special case where $\thickmuskip=2mu \medmuskip=2mu r=1$, the block-diagonal structure simplifies to a conventional unconstrained orthogonal matrix. In contrast, approaches like Stable Diffusion~\cite{rombach2022high} and DALL-E2~\cite{ramesh2022hierarchical} operate in latent spaces. These findings suggest that neural directionality predominantly encodes the semantic content of images, implying that adjusting neuron angles is likely a particularly effective strategy for manipulating image semantics. For enhanced computational practicality, we utilize the Cayley parameterization as an efficient means to generate orthogonal matrices. For classification tasks, redundancy-free neurons are important, but arranging them to be maximally spaced (as minimum hyperspherical energy would dictate) is not the right objective during finetuning, since this would compromise and potentially eliminate valuable pretraining structure. 
\begin{abstract}
Large-scale text-to-image diffusion models have demonstrated remarkable proficiency in synthesizing photorealistic visuals from textual descriptions. The COFT approach thus operates through two explicit regularizations: maintaining a minimal hyperspherical energy discrepancy and constraining deviation from the original pretrained parameters. Unlike these stable convergence trajectories, ControlNet demonstrates an abrupt improvement in controllability at epoch 8, consistent with the sudden convergence effect described in \cite{zhang2023adding}. To address this, we leverage an important invariance property of hyperspherical energy: when an identical orthogonal transformation is applied to all neurons, the pairwise hyperspherical similarities are theoretically conserved. We next examine the parameter efficiency of OFT in comparison with LoRA. Adaptation-based finetuning strategies (\emph{e.g.}, \cite{hulora2022,houlsby2019parameter,pfeiffer2020adapterfusion}) have been extensively explored in natural language processing. In the following, we detail approaches for preserving orthogonality in a way that is computationally tractable. Drawing a parallel with LoRA’s use of zero initialization, it is necessary for OFT to ensure that fine-tuning commences precisely from the pretrained $\bm{W}^0$. Eq.~\eqref{eq:oft_general}), in which only $\bm{R}$ was optimized, here both $\bm{R}$ and $\bm{D}$ are updated during learning. This modification is based on the observation that scaling neurons does not impact hyperspherical energy, since any change in magnitude is neutralized through normalization. As an evaluation criterion, we calculate the average $\ell_2$ distance between the control and predicted face landmarks. Despite these advancements, fully preserving the fine details that characterize a subject’s identity remains a challenge. More concretely, the orthogonal matrix is parameterized as $\thickmuskip=2mu \medmuskip=2mu \bm{R}=(\bm{I}+\bm{Q})(I-\bm{Q})^{-1}$, with $\bm{Q}$ being skew-symmetric, i.e., $\thickmuskip=2mu \medmuskip=2mu \bm{Q}=-\bm{Q}^\top$. Identifying strategies to further enhance parameter efficiency with reduced bias is an important challenge moving forward. This evaluation confirms that the OFT framework offers superior convergence and stability, while consistently producing higher-quality final results. Unlike the original OFT setup (cf. To control for stochasticity, we repeat finetuning of the same pretrained model 30 times with distinct random seeds for every approach. To provide a clearer, more tangible illustration of the advantages brought by OFT, we display several randomly selected samples for subject-driven generation in Figure~\ref{fig:dreambooth_final}. To maintain orthogonality for $\bm{R}$, differential orthogonalization techniques, as reviewed in \cite{liu2021orthogonal,lezcano2019cheap}, are applicable. The OFT framework is, by design, applicable to any neural network structure. Successful finetuning thus hinges on balancing adaptability with preservation of learned structure; our OFT paradigm is constructed precisely to strike this balance. Conversely, if all blocks in $\bm{R}$ are identical, the same orthogonal transformation is applied across all channels. Moreover, the block-diagonal approach can be interpreted as imposing stricter regularization on the orthogonal transformation matrix. While many standard finetuning approaches implicitly adopt the second form of constraint, COFT introduces it as an explicit design choice. Across all three control scenarios, OFT and COFT demonstrate a qualitative advantage over baseline models, underscoring the efficacy of the OFT framework for controllable image generation. Despite its benefits, OFT raises several noteworthy open research questions. These observations support that the OFT framework attains rapid and reliable convergence when tackling subject-driven generation tasks. In Constrained Orthogonal Finetuning (COFT), the forward computation proceeds as:

\begin{equation}\label{eq:coft_general}
    \footnotesize
    \bm{z}=\bm{W}^\top\bm{x}=(\bm{R}\cdot\bm{W}^0)^\top\bm{x},~~~\text{s.t.}~\bm{R}^\top\bm{R}=\bm{R}\bm{R}^\top=\bm{I},~\norm{\bm{R}-\bm{I}}\leq \epsilon
\end{equation}

where both an orthogonality constraint and a proximity constraint (enforced via an $\epsilon$-bound to the identity) are imposed. We partially mitigate this by adopting a block diagonal parametrization, but finding a more differentiable and faster solution to the matrix inversion problem remains an open issue. \item For evaluating controllable generation, we propose a novel metric for control consistency, which quantifies the degree of controllability. Lastly, the current parameter efficiency of OFT depends heavily on the block diagonal design, which can inherently add bias and restrict flexibility. \textbf{Convergence}. By contrast, omitting the orthogonality requirement from OFT results in a complete inability to generate plausible images or achieve any controllability. Contrastingly, OFT applies a different form of regularization, namely constraining the similarity of neuron pairs through their hyperspherical energy: $\thickmuskip=2mu \medmuskip=2mu \| \text{HE}(\bm{M})-\text{HE}(\bm{M}^0) \|=0$, where $\text{HE}(\cdot)$ indicates the hyperspherical energy of the corresponding weight matrix. \textbf{Why OFT converges faster}. The main contributions of our work are summarized as follows:

\begin{itemize}[leftmargin=*,nosep]

    \item We introduce Orthogonal Finetuning, a new method designed to enhance the controllability of text-to-image diffusion models. For LoRA, given a low-rank parameter $r'$, its trainable parameter count is $\thickmuskip=2mu \medmuskip=2mu r'(d+n)$. \item \textbf{Controllable generation}~\cite{zhang2023adding,mou2023t2i}: With access to an external control signal (\emph{e.g.}, edge maps, segmentation annotations), the goal here is to synthesize images that not only adhere to the input signal but also conform to a text prompt. Similarly, T2I-Adapter~\cite{mou2023t2i} refines pretrained diffusion models to improve controllable image synthesis. It is worth highlighting that OFT and COFT utilize a parameter count comparable to LoRA (and substantially less than ControlNet), yet attain markedly improved convergence efficiency relative to prior methods. Visual results in Figure~\ref{fig:face_convergence_qual} indicate that both OFT and COFT exhibit stable convergence, with synthesized face poses incrementally aligning with the control landmarks. To ensure initialization of the orthogonal matrix $\bm{R}$ as the identity, we simply set $\bm{Q}$ to the zero matrix at initialization. By contrast, our work targets the finetuning of text-to-image diffusion models, prioritizing both precise controllability and robust generative performance on downstream tasks. One core challenge stems from the absence of an objective metric to assess how well the pretrained generative proficiency is preserved after finetuning. Inversion-based techniques~\cite{choi2021ilvr,dhariwal2021diffusion,ramesh2022hierarchical,gal2022image} adjust the contextual elements of an image while retaining the subject. Although OFT already exhibits robust performance, we find that COFT delivers improved finetuning stability owing to the added explicit regularization of deviation. We underline that our approach to quantitative analysis of controllable generation is among the paper’s key innovations, as it precisely reflects the control fidelity of finetuned models, bounded only by the performance of the underlying segmentation or detection system used for evaluation. Additional visual examples are available in Appendix~\ref{app:more_qualitative}. It is evident that COFT and OFT exhibit notably stable finetuning behavior for the diffusion model. Motivated by empirical findings that hyperspherical similarity effectively represents semantic information~\cite{liu2017deep,liu2018decoupled,chen2020angular}, we employ the notion of hyperspherical energy~\cite{liu2018learning} to capture the pairwise relationships among neurons. \textbf{Quantitative comparison}. \end{abstract}

\section{Introduction}

Recent advancements in text-to-image diffusion models~\cite{saharia2022photorealistic,ramesh2022hierarchical,rombach2022high} have established state-of-the-art results in generating high-quality images conditioned on text prompts. Nevertheless, even with the Cayley transform, using a large dimension $d$ may render $\bm{R}$ highly parameter inefficient. Further guidance comes from the literature on orthogonal over-parameterized training~\cite{liu2021orthogonal}, where classification networks are trained from random initialization using orthogonal transformations. Notably, OFT introduces no extra computational expense during inference at test time. These findings underscore the necessity of the orthogonality constraint within OFT. To illustrate, consider a standard fully connected layer with weights $\thickmuskip=2mu \medmuskip=2mu \bm{W}=\{\bm{w}_1,\cdots,\bm{w}_n\}\in\mathbb{R}^{d\times n}$, where each $\bm{w}_i\in\mathbb{R}^{d}$ stands for the $i$-th neuron and $\bm{W}^0$ is the pretrained counterpart. In contrast to ControlNet, the OFT framework does not increase computational demand during inference. \subsection{Controllable Generation}

\textbf{Settings}. First, as depicted in Figure~\ref{fig:toy_ae}, the most substantial changes in the model’s representations arise from altering the directions of neurons—an adaptation uniquely facilitated by OFT. To address this, we propose a rigorous adaptation technique called Orthogonal Finetuning~(OFT) that enables text-to-image diffusion models to be fine-tuned for specific tasks. In particular, we examine two major categories of finetuning tasks for text-to-image models: subject-driven generation, aimed at producing images of a given subject in new scenarios from a handful of example images and a prompt, and controllable generation, which involves incorporating auxiliary control signals. The principal mechanism is to learn an orthogonal transformation that is consistent across a layer, thereby preserving the angular relationships between neurons. For COFT, an appropriate choice of $\epsilon$ further simplifies the selection of both learning rate and iteration steps. When each block in $\bm{R}$ corresponds to a single input channel, every neuron channel is individually transformed, echoing the functionality of depth-wise convolution kernels~\cite{chollet2017xception}. Empirical support for this perspective can also be found in \cite{liu2021orthogonal}. Firstly, while~\cite{liu2021orthogonal} focuses on reducing hyperspherical energy, OFT is designed to preserve the pretrained model’s hyperspherical energy, thereby safeguarding the intrinsic structure established during pretraining—a critical consideration, as minimizing energy in diffusion model finetuning may disrupt essential semantic configurations. This alternative formulation permits direct enforcement of the constraint on $\bm{Q}$ via projected gradient descent. Differing from this, OFT introduces a layer-shared orthogonal transformation, modifying neuron weights multiplicatively. The outcomes shown in Figure~\ref{fig:toy_ae} reveal that reconstructing images with only angular components of neuron activations almost perfectly recovers the originals, whereas the magnitude holds negligible relevant information. Pivotal Tuning~\cite{roich2022pivotal} addresses this by finetuning a generator around an inverted latent code, supplemented with an identity-preserving regularization term. During model deployment, one can simply left-multiply the optimized orthogonal matrix $\bm{R}$ with the original weight matrix $\bm{W}^0$, resulting in an equivalent parameter matrix $\thickmuskip=2mu \medmuskip=2mu \bm{W}=\bm{R}\bm{W}^0$. Our empirical analysis shows that this property is vital for retaining the generative semantics of the original model. For diversity, the LPIPS metric evaluates average pairwise cosine similarity among outputs generated for the same subject and prompt. Although~\cite{casanova2021instance} can model variations of a given instance, it may compromise specificity. For evaluation, we test three distinct strategies for feature map generation: (a) the standard inner product as used during training, (b) reliance solely on the magnitude of the outputs, and (c) use of purely angular information. \textbf{Zero inference-time overhead}. In the context of Transformers, LoRA is typically integrated into attention layers~\cite{hulora2022}. First, OFT enforces orthogonality using Cayley parametrization, which necessitates a matrix inversion, marginally constraining the scalability of the method. To overcome these limitations, our study focuses on two particular types of text-to-image generation tasks:

\begin{itemize}[leftmargin=*,nosep]

    \item \textbf{Subject-driven generation}~\cite{ruiz2023dreambooth}: Provided with a small collection of images featuring a specific subject, this task involves generating images of the same subject situated in varied environments, guided by a corresponding text description. For our experiments, Stable Diffusion v1.5~\cite{rombach2022high} serves as the foundational text-to-image generation model. Our approach is applied to two specific text-to-image tasks: subject-driven generation and controllable generation. \subsection{Subject-driven Generation}

\textbf{Configuration}. To provide a fair comparison with LoRA, the OFT and COFT methods are exclusively applied to the identical network layers as those used by LoRA. Within this context, OFT seeks to minimize the discrepancy in hyperspherical energy between the adjusted and original parameters as follows:

\begin{equation}\label{eq:oft_principle}
\footnotesize
\min_{\bm{W}} \left\| \text{HE}(\bm{W})-\text{HE}(\bm{W}^0)\right\|~~\Leftrightarrow~~\min_{\bm{W}} \bigg{\|} \sum_{i\neq j} \|\hat{\bm{w}}_i-\hat{\bm{w}}_j\|^{-1}-\sum_{i\neq j} \|\hat{\bm{w}}^0_i-\hat{\bm{w}}^0_j\|^{-1}\bigg{\|}
\end{equation}

Here, the normalized neuron is written as $\thickmuskip=2mu \medmuskip=2mu \hat{\bm{w}}_i:=\bm{w}_i/\|\bm{w}_i\|$, and the hyperspherical energy for a weight matrix $\bm{W}$ is given by $\thickmuskip=2mu \medmuskip=2mu \text{HE}(\bm{W}):= \sum_{i\neq j} \|\hat{\bm{w}}_i-\hat{\bm{w}}_j\|^{-1}$. \textbf{Model finetuning}. We propose a metric for control consistency to systematically assess controllable generation capabilities.